% =====================================================================
% 12a_livrables.tex — livrables et organisation du dépôt
% =====================================================================
\chapter{Livrables et organisation du dépôt}
\label{ann:livrables}

\section{Arborescence}

Le dépôt est organisé par partie, avec un script maître et des sorties de
preuve. Les scripts SQL sont à la racine ; le rapport LaTeX est dans le
sous-dossier \code{Rapport finale}.

\begin{lstlisting}[caption={Organisation des fichiers du projet.},label={lst:arborescence},language={}]
D:\bdda\
|-- 00_creation_utilisateur.sql
|-- partie_A_schema.sql
|-- partie_B_manipulations.sql
|-- partie_C_requetes.sql
|-- partie_D_programmation.sql
|-- partie_E_objet_avance.sql
|-- projet_bdda_banque.sql
|-- demo_soutenance.sql
|-- purge_donnees.sql
|-- reset_schema.sql
|-- GUIDE_SOUTENANCE.md
|-- CAPTURES_A_FAIRE.md
|-- ETAPE_0_CADRAGE.md
|-- sujet.txt
|-- uml\
|   |-- banque.puml
|   |-- banque.mermaid
|   `-- banque.txt
|-- _travail\
|   |-- _A.out ... _E.out
|   |-- _FULL.out
|   |-- _DEMO.out
|   `-- _runFULL.bat
`-- Rapport finale\
    |-- main.tex
    |-- config\
    |-- chapters\
    |-- images\
    |-- bibliography\
    `-- README.md
\end{lstlisting}

\section{Description des livrables}

{\small
\setlength{\tabcolsep}{3pt}
\begin{longtable}{@{}L{5.9cm}L{2.9cm}L{6.1cm}@{}}
\caption{Description des fichiers livrés.}\label{tab:livrables}\\
\toprule
\entete{Fichier} & \entete{Rôle} & \entete{Contenu} \\
\midrule
\endfirsthead
\toprule
\entete{Fichier} & \entete{Rôle} & \entete{Contenu} \\
\midrule
\endhead
\bottomrule
\endfoot
\code{00\_creation\_utilisateur.sql} & Administration &
Création idempotente de \code{MASTER1}, privilèges et quota. À exécuter par un
administrateur, non exécuté pendant le projet. \\
\code{partie\_A\_schema.sql} & Schéma &
Types, tables, contraintes, séquences et jeu de données. \\
\code{partie\_B\_manipulations.sql} & Manipulations &
B5a à B5f, ajout du client 99, contrainte différée, \code{VARRAY}, suppression
en cascade. \\
\code{partie\_C\_requetes.sql} & Requêtes &
25 requêtes SQL, dont les variantes 6.3a, 6.7b, 6.9a, 6.11b, 6.16b, 6.18b,
6.19c, 6.21b. \\
\code{partie\_D\_programmation.sql} & Programmation &
Sous-programmes, blocs anonymes, méthodes, triggers et tests D9. \\
\code{partie\_E\_objet\_avance.sql} & Objet avancé &
E10, E11, E12 : parrainage, relevés multi-niveaux, vue objet et mutation. \\
\code{projet\_bdda\_banque.sql} & Script maître &
Concaténation A à E, exécutable seul, avec contrôle final. \\
\code{demo\_soutenance.sql} & Démonstration &
Séquence de dix minutes, mouvements annulés par \code{ROLLBACK}. \\
\code{purge\_donnees.sql} & Maintenance &
Vidage des données sans suppression du schéma. \\
\code{reset\_schema.sql} & Maintenance &
Remise à zéro du schéma pour un rejeu propre. \\
\code{GUIDE\_SOUTENANCE.md} & Guide &
Démonstration, 45 questions-réponses, pièges et chiffres clés. \\
\code{CAPTURES\_A\_FAIRE.md} & Preuves &
Liste ordonnée de 66 captures, priorités et chronologie. \\
\code{ETAPE\_0\_CADRAGE.md} & Conception &
Hypothèses H1 à H14, choix D1 à D20, environnement vérifié. \\
\code{uml/} & Modèle &
Diagramme PlantUML, version Mermaid et description textuelle. \\
\code{\_travail/} & Preuves &
Sorties \code{\_A.out} à \code{\_E.out}, \code{\_FULL.out},
\code{\_DEMO.out}. \\
\code{Rapport finale/} & Rapport &
Sources LaTeX, figures, bibliographie, PDF compilé et README. \\
\end{longtable}
}

\section{Environnement d'exécution}

L'environnement cible est Oracle Database 21c Express Edition, avec la base
pluggable \code{XEPDB1} et l'utilisateur \code{MASTER1}. La connexion utilisée
pendant le projet est :

\begin{lstlisting}[caption={Connexion SQL*Plus au schéma du projet.},label={lst:connexion},language={}]
sqlplus MASTER1/system@localhost:1521/XEPDB1
\end{lstlisting}

\begin{attentionbox}[Ne pas utiliser SYSTEM]
Les scripts applicatifs doivent être exécutés avec \code{MASTER1}. Une
connexion \code{SYSTEM} ou \code{SYS} créerait les objets dans le mauvais
schéma et fausserait les contrôles. Seul
\code{00\_creation\_utilisateur.sql} est destiné à un administrateur.
\end{attentionbox}

Les scripts commencent par les paramètres SQL*Plus suivants :

\begin{lstlisting}[caption={Paramètres communs aux scripts SQL.},label={lst:parametres-sql}]
SET SERVEROUTPUT ON SIZE UNLIMITED
SET LINESIZE 200
SET PAGESIZE 100
SET DEFINE OFF
SET FEEDBACK ON
SET VERIFY OFF
\end{lstlisting}

\code{SERVEROUTPUT} active \code{DBMS\_OUTPUT}, \code{LINESIZE} évite les
retours à la ligne parasites, \code{DEFINE OFF} désactive la substitution
\code{\&}, et \code{VERIFY OFF} allège les sorties.

\section{Rejeu complet}

Le rejeu complet se fait avec le script maître. Il supprime puis recrée les
objets, insère les données et exécute les parties B à E. Le contrôle final
vérifie l'absence d'objets invalides.

\begin{lstlisting}[caption={Rejeu complet et contrôle final.},label={lst:rejeu},language={}]
sqlplus -S -L MASTER1/system@localhost:1521/XEPDB1 @projet_bdda_banque.sql

SELECT COUNT(*) AS objets_invalides
  FROM user_objects
 WHERE status <> 'VALID';
\end{lstlisting}

Le script peut aussi être exécuté partie par partie, dans l'ordre A, B, C, D,
E. Cet ordre est important : la partie C doit être capturée avant les tests
destructifs de la partie D. Le script \code{reset\_schema.sql} remet le schéma
à zéro si nécessaire.

\section{Preuves conservées}

Les sorties complètes sont conservées dans \code{\_travail}. Elles permettent
de montrer un résultat sans modifier l'état courant de la base.

\begin{table}[htbp]
  \centering
  \small
  \caption{Sorties de preuve conservées.}
  \label{tab:preuves}
  \begin{tabularx}{\textwidth}{@{}L{3.8cm}L{3.2cm}X@{}}
    \toprule
    \entete{Fichier} & \entete{Portée} & \entete{Utilisation} \\
    \midrule
    \code{\_A.out} & Partie A & Création du schéma et données. \\
    \code{\_B.out} & Partie B & B5a à B5f. \\
    \code{\_C.out} & Partie C & Requêtes 6.1 à 6.25, avant D. \\
    \code{\_D.out} & Partie D & Méthodes, triggers, clôture. \\
    \code{\_E.out} & Partie E & Parrainage, relevés, vue objet. \\
    \code{\_FULL.out} & A à E & Rejeu complet, 56\,549 octets. \\
    \code{\_DEMO.out} & Démonstration & Sortie de \code{demo\_soutenance.sql}. \\
    \code{\_runFULL.bat} & Rejeu & Lancement automatisé du script maître. \\
    \code{\_RESET.out} & Maintenance & Sortie de la remise à zéro. \\
    \code{\_verif\_reset.out} & Maintenance & Contrôle après remise à zéro. \\
    \bottomrule
  \end{tabularx}
\end{table}

\section{Compilation du rapport}

Le rapport utilise XeLaTeX et Biber. La commande recommandée est :

\begin{lstlisting}[caption={Compilation du rapport.},label={lst:compilation},language={}]
latexmk -xelatex -interaction=nonstopmode -halt-on-error main.tex
\end{lstlisting}

En cas de compilation manuelle, l'ordre est : \code{xelatex}, \code{biber},
\code{xelatex}, \code{xelatex}. Le PDF produit est \code{main.pdf}.
